\section{External Behavioral Validation on CLUTRR}
\label{sec:clutrr-external-validation}

The preceding audits establish that representation-level holdouts can fail to
exclude familiar semantics, but protocol exposure alone does not show that a
learner uses that exposure.  We therefore ran a controlled intervention on
CLUTRR \cite{sinha2019clutrr}, a public benchmark of kinship reasoning from
generated text.  The intervention asks a narrow behavioral question:
\emph{does training on one set of surface realizations of a relation pattern
improve accuracy on source-group-disjoint realizations of the same semantic
pattern?}  We report a preregistered pilot (V3) and a subsequent,
target-disjoint directional confirmation (V4).  The studies are not pooled.

\subsection{Task, Semantic Unit, and Intervention}

We freeze CLUTRR Task~1 at relation length three, without distractor facts.
Each example contains a three-edge kinship chain and a query relation drawn
from 23 classes.  Entity names are canonicalized by order of appearance.  The
semantic identity is the ordered tuple consisting of the target relation and
the three edge relations; sentence order is therefore a surface property, not
part of the registered semantic identity.

For every three-sentence source story, we enumerate all six sentence
permutations.  The six permutations share a source-group identifier.  A target
pattern is eligible only when it has at least 45 complete source groups.  Its
evaluation set contains 200 rows from the first 34 sorted groups: all six
permutations from groups 0--32 and two permutations from group~33.  Its
training exposure contains 64 rows from the next 11 groups: all six
permutations from groups 34--43 and four from group~44.  Evaluation and
exposure source groups are consequently disjoint.

Targets are paired only within the same output relation.  For each model
seed, a frozen SHA-256 bit exposes one pair member in Arm~A and the other in
Arm~B; the assignment is reversed for the complementary member.  Both arms
share a 22,720-row base containing no target semantics and append 1,280
target-exposure rows, giving 24,000 training rows per run with identical
label counts.  Initialization, all 30 epoch permutations, tensor shapes, and
dropout RNG trajectories are paired across complementary runs.

We train a fixed four-layer byte Transformer with model dimension 128, four
attention heads, feed-forward dimension 512, and 0.1 dropout.  All runs use
AdamW, batch size 128, learning rate $3\times10^{-4}$, deterministic
float32 CPU execution, and the fixed epoch-30 checkpoint after 5,640 updates.
No validation selection or early stopping is used.  The evaluation manifest
is identical across arms within a block.

\subsection{Estimand and Registered Inference}

Let $A_{s,f}(e)$ be row-weighted accuracy for model seed $s$, target semantic
pattern $f$, and exposure status $e\in\{0,1\}$.  For $S$ frozen seeds and
$F$ frozen targets, the sole primary estimand is
\begin{equation}
\tau
=
\frac{100}{SF}
\sum_{s=1}^{S}\sum_{f=1}^{F}
\left[A_{s,f}(1)-A_{s,f}(0)\right]
\quad\text{percentage points}.
\label{eq:clutrr-exposure-estimand}
\end{equation}
The 95\% interval uses 20,000 hierarchical percentile-bootstrap replicates.
Each replicate resamples model seeds, complete same-label target pairs, and
source groups within each selected pair member.  Source-group draws are shared
between exposed and unexposed predictions and across selected seed
occurrences; all rows in a sampled group occurrence are retained.

\begin{table}[t]
\centering
\small
\setlength{\tabcolsep}{4pt}
\begin{tabular}{lrrrrl}
\toprule
Study & Pairs & Targets & Runs & Effect (pp) & 95\% CI / Gate \\
\midrule
V3 pilot & 20 & 40 & 16 & $+15.338$ &
$[9.877,21.369]$; precision No-Go \\
V4 confirm. & 40 & 80 & 32 & $+23.339$ &
$[17.750,29.226]$; directional Go \\
\bottomrule
\end{tabular}
\caption{CLUTRR exposure interventions.  V3 and V4 use disjoint target
semantic patterns and are not pooled.}
\label{tab:clutrr-external-validation}
\end{table}

\subsection{V3: Positive Direction, Registered Precision No-Go}

V3 froze 20 same-label pairs, 40 targets, eight model seeds, and two arms,
for 16 runs.  Its preregistration admitted positive, negative, and
equivalence conclusions only after a precision requirement: the 95\%
interval width had to be at most two points.  The observed effect was
$+15.3375$ points with 95\% CI $[9.8769,21.3693]$.  Although the entire
interval was positive, its width was 11.4924 points, so the registered
classification was \texttt{imprecise\_no\_go}.  We preserve that decision
rather than relabeling the pilot after seeing a large estimate.

After V3 was closed, an explicitly exploratory variance diagnosis attributed
94.40\% of bootstrap variance to target-pair resampling, 4.90\% to model-seed
resampling, and 0.70\% to source-group resampling.  A calibrated prospective
projection showed that reproducing the original two-point precision
requirement would need at least 634 independent pairs even with infinitely
many seeds, while only 139 V3-disjoint eligible pairs remained.  Using all
remaining pairs would still project to a 4.269-point interval width.  This
analysis closed the strict-precision replication route; it did not change the
V3 result.

\subsection{V4: Prospective Target-Disjoint Confirmation}

We used V3 only as a declared pilot to freeze a new directional design before
any V4 training outcome existed.  The design effect was the smaller of half
the V3 point estimate and the V3 lower interval endpoint, or 7.66875 points.
A 100,000-draw Monte Carlo Gate compared one to four V3-sized blocks.  One
block (20 pairs) had estimated power 0.7440 and failed.  Two blocks (40 pairs)
had power 0.9583 with 95\% Wilson lower bound 0.9570 and were selected by the
registered minimum-passing-block rule.  All 80 V4 target patterns were
semantically disjoint from V3 targets.  The V4 primary Gate required
$\tau\geq1$ point and a two-sided 95\% interval lower endpoint above zero.

All 32 V4 runs reached the fixed checkpoint and passed post-training
integrity checks for model and optimizer finiteness, frozen epoch orders,
four-way paired RNG trajectories, prediction ordering, independent metric
recomputation, and artifact hashes.  The confirmatory effect was
$+23.3391$ points with 95\% CI $[17.7498,29.2257]$, satisfying the registered
\texttt{directional\_positive\_go} rule.  Descriptively, both block estimates
were positive ($+25.0109$ and $+21.6672$ points), as were all eight seed
estimates and all 40 pair estimates; 76 of 80 individual target-pattern
estimates were positive.  These sign counts were not additional
multiplicity-free tests and did not enter the Gate.

The secondary proper losses move in the same favorable direction:
exposed-minus-unexposed cross entropy is $-1.4938$ and multiclass Brier score
is $-0.4224$.  ECE-15 changes by $-0.2018$.  These metrics were registered as
descriptive and cannot alter the accuracy-based primary conclusion.

\subsection{Interpretation and Scope}

The intervention identifies an endpoint effect of semantic exposure in the
frozen CLUTRR setting.  A model trained on one collection of stories
instantiating a relation composition performs better on different source
groups and sentence permutations instantiating that same composition.  The
result therefore demonstrates behavioral use of semantic familiarity beyond
exact-string reuse.  It also supplies the external public-benchmark behavior
that the earlier structural audits lacked.

The claim remains deliberately bounded.  The experiment covers one CLUTRR
task, relation length, model architecture, exposure dose, and optimization
budget.  It does not show that every form of semantic overlap improves
accuracy, that the model stores a literal lookup table, or that the same
effect holds for unseen relation lengths, family-disjoint reasoning, other
architectures, or natural text.  In particular, the near-zero registered
Boolean intervention remains a counterexample to the universal shortcut
``overlap implies memorization.''  The combined evidence supports the weaker
and more useful conclusion: semantic exposure can have a large causal
behavioral effect, so benchmark novelty and behavioral exploitation must be
measured separately.
